%% THERMODYNAMIQUE %%
%
%% Dans ce fichier, les �ventuels accents sont cod�s
%% en ISO-8859-15 (aka latin9)
%% La ligne suivante donne trois exemples 
%% e accent aigu : �, e accent grave : �, e-dans-l'o : �
\section{Temp�rature statique (lorsque ventilation suffisante)}

\itemstyle{Op\'erateur~:} \operfmt{Tair\_avion}
\index{Tair\_avion}

%%%%%OBJET%%%%
\itemstyle{Objet~:}
Calcul de la temp�rature statique de l'air, � 
partir de la mesure de temp�rature totale et de la surpression dynamique.\newline
Ce calcul s'applique aux capteurs de temp�rature par contact
de type Rosemount par exemple.\newline
Le calcul n'est fait que lorsque la surpression dynamique est
sup�rieure � 6 hPa. Dans le cas contraire, signe d'une
ventilation insuffisante du capteur, la valeur retourn�e est
le code "valeur manquante".


%%%%%ARGUMENTS%%%%
\bigskip
%% Dans ce fichier, les �ventuels accents sont cod�s
%% en ISO-8859-15 (aka latin9)
%% La ligne suivante donne trois exemples 
%% e accent aigu : �, e accent grave : �, e-dans-l'o : �
%%%%%ENTREES%%%%
\begin{tabarg}
\multicolumn{3}{l}{\titretab {Entr\'ees}}\\
\hline
$T_t$ & Vecteur, \texttt{float} &  Temp�rature totale ($K$ ou $^\circ C$)\\
$\Delta P$ & Vecteur, \texttt{float} & Pression dynamique ($hPa$)\\
$P_s$ & Vecteur, \texttt{float} & Pression statique ($hPa$)\\
$r_f$ & Coeff., \texttt{float} & Coefficient de r�cup�ration du capteur utilis�\\
$R_a/c_{pa}$ & Coeff., \texttt{float} & Rapport de la constante sp�cifique de
l'air sec sur la chaleur massique � pression constante de l'air sec\\
\end{tabarg}

\medskip
%%%%%SORTIES%%%%
\begin{tabarg}
\multicolumn{3}{l}{\titretab {Sorties}}\\
\hline
$T_s$ & Vecteur, \texttt{float} & Temp�rature statique de l'air (m�me unit� que $T_t$)\\
\end{tabarg}


%%%%%FORMULE%%%%
\itemstyle{Formule}

%% Formule de l'operateur : Tair_avion
%% Dans ce fichier, les �ventuels accents sont cod�s
%% en ISO-8859-15 (aka latin9)
%% La ligne suivante donne trois exemples 
%% e accent aigu : �, e accent grave : �, e-dans-l'o : �
Si besoin est, $T_t$ est converti en Kelvin, puis
\[
T_s =
\begin{cases}
\frac{T_t}
     {1 + r_f \left( \left(1 + \frac{\Delta P}{P_s}\right)
        ^{R_a/c_{pa}} - 1\right)} \qquad \text{ si } \Delta P > 6 \\
\text{valeur\_manquante} \qquad \text{ sinon } \\
\end{cases}
\]


%%%%%AUTEUR%%%%
\itemstyle{Auteur~:}
CNRM/GMEI/TRAMM

%%%%%REFERENCE%%%%
\itemstyle{R\'ef\'erences~:}
